\documentclass[11pt,a4paper]{article}

% --------------------------------------------------
% Encodage et langue
% --------------------------------------------------
\usepackage[utf8]{inputenc}
\usepackage[T1]{fontenc}
\usepackage[french]{babel}

% --------------------------------------------------
% Mise en page
% --------------------------------------------------
\usepackage[a4paper,margin=2cm]{geometry}
\usepackage{lmodern}
\usepackage{microtype}
\usepackage{setspace}
\setstretch{1.08}

% --------------------------------------------------
% Maths et physique
% --------------------------------------------------
\usepackage{amsmath,amssymb,amsthm,mathtools}
\usepackage{siunitx}
\sisetup{locale=FR}
\usepackage{physics}

% --------------------------------------------------
% Tableaux, listes, couleurs
% --------------------------------------------------
\usepackage{array}
\usepackage{tabularx}
\usepackage{booktabs}
\usepackage{enumitem}
\usepackage{xcolor}
\usepackage{tcolorbox}
\tcbuselibrary{breakable,skins}
\usepackage{multicol}

% --------------------------------------------------
% Figures et TikZ
% --------------------------------------------------
\usepackage{graphicx}
\usepackage{tikz}
\usetikzlibrary{arrows.meta,calc,decorations.pathmorphing,patterns,positioning}
\usepackage{pgfplots}
\pgfplotsset{compat=1.18}

% --------------------------------------------------
% Hyperliens et en-têtes
% --------------------------------------------------
\usepackage{hyperref}
\hypersetup{
    colorlinks=true,
    linkcolor=blue!60!black,
    urlcolor=blue!60!black,
    citecolor=blue!60!black
}

\usepackage{fancyhdr}
\usepackage{lastpage}
\setlength{\headheight}{14pt}
\pagestyle{fancy}
\fancyhf{}
\fancyhead[L]{Première générale -- Spécialité Physique-Chimie}
\fancyhead[C]{Cours complet -- Mouvement, interactions et énergie}
\fancyhead[R]{Page \thepage/\pageref{LastPage}}
\fancyfoot[C]{\textit{Cours détaillé, démonstrations, méthodes, exemples et exercices}}

% --------------------------------------------------
% Théorèmes et styles
% --------------------------------------------------
\newtheorem{definition}{Définition}[section]
\newtheorem{propriete}{Propriété}[section]
\newtheorem{theoreme}{Théorème}[section]
\newtheorem{remarque}{Remarque}[section]
\newtheorem{exemple}{Exemple}[section]

\newtcolorbox{objectifbox}{
    breakable,
    enhanced,
    colback=blue!5,
    colframe=blue!60!black,
    title=Objectifs du chapitre,
    fonttitle=\bfseries
}

\newtcolorbox{methodebox}{
    breakable,
    enhanced,
    colback=green!5,
    colframe=green!45!black,
    title=Méthode,
    fonttitle=\bfseries
}

\newtcolorbox{attentionbox}{
    breakable,
    enhanced,
    colback=red!5,
    colframe=red!60!black,
    title=Attention,
    fonttitle=\bfseries
}

\newtcolorbox{terminalebox}{
    breakable,
    enhanced,
    colback=violet!5,
    colframe=violet!60!black,
    title=Ouverture vers la Terminale,
    fonttitle=\bfseries
}

\newtcolorbox{applicationbox}{
    breakable,
    enhanced,
    colback=orange!5,
    colframe=orange!70!black,
    title=Application,
    fonttitle=\bfseries
}

\newcommand{\R}{\mathbb{R}}
\newcommand{\Ec}{E_\mathrm{c}}
\newcommand{\Ep}{E_\mathrm{p}}
\newcommand{\Em}{E_\mathrm{m}}

\title{\Huge Cours très détaillé\\[0.3em]
\LARGE Mouvement et interactions\\
\Large avec les aspects énergétiques des phénomènes mécaniques et électriques}
\author{}
\date{}

\begin{document}

\maketitle
\tableofcontents
\newpage

% ==================================================
\section*{Introduction générale}
% ==================================================

La mécanique est une partie centrale de la physique. Elle cherche à répondre à des questions simples en apparence :

\begin{itemize}
    \item pourquoi un objet reste-t-il immobile ou se met-il en mouvement ?
    \item pourquoi sa vitesse change-t-elle ?
    \item comment modéliser l'action d'un autre objet ou d'un milieu sur lui ?
    \item comment relier les forces, le mouvement et l'énergie ?
\end{itemize}

Ce thème est fondamental car il introduit des idées qui seront développées en Terminale : lois de Newton sous forme plus complète, champs, travail, énergie, mouvement dans des situations plus riches.

\begin{objectifbox}
À la fin de ce chapitre, il faut être capable de :
\begin{itemize}
    \item distinguer interaction, force et champ ;
    \item utiliser les lois de gravitation et d'interaction électrostatique ;
    \item comprendre la description d'un fluide au repos ;
    \item exploiter la relation entre variation de vitesse et somme des forces ;
    \item utiliser les notions d'énergie cinétique, de travail, d'énergie potentielle et d'énergie mécanique ;
    \item faire un lien clair entre description du mouvement et bilan énergétique ;
    \item éviter les erreurs classiques de raisonnement spontané en mécanique.
\end{itemize}
\end{objectifbox}

% ==================================================
\section{Rappels de seconde : référentiel, position, vitesse, forces}
% ==================================================

Avant d'aller plus loin, il faut bien fixer le vocabulaire.

\begin{definition}
Un \textbf{référentiel} est un solide ou un ensemble de points par rapport auquel on décrit le mouvement.
\end{definition}

En pratique, au lycée, on travaille souvent dans le référentiel terrestre, supposé galiléen dans de nombreuses situations usuelles.

\begin{definition}
Le \textbf{vecteur position} d'un point matériel donne sa position dans l'espace à un instant donné.
\end{definition}

\begin{definition}
Le \textbf{vecteur vitesse} indique :
\begin{itemize}
    \item la direction du mouvement ;
    \item le sens du mouvement ;
    \item la valeur de la vitesse.
\end{itemize}
\end{definition}

\begin{definition}
Une \textbf{force} modélise l'action d'un objet ou d'un milieu sur un système.
\end{definition}

\begin{remarque}
Une force n'est pas « la cause magique du mouvement ». Plus précisément, une force modifie l'état du mouvement, c'est-à-dire la vitesse.  
Un objet peut se déplacer sans force résultante non nulle : il peut alors conserver sa vitesse.
\end{remarque}

\begin{attentionbox}
Erreur fréquente : penser que « s'il y a mouvement, alors il y a forcément une force dans le sens du mouvement ».  
C'est faux. Une force est nécessaire pour \textbf{modifier} la vitesse, pas simplement pour « entretenir » le mouvement.
\end{attentionbox}

% ==================================================
\section{Interactions fondamentales et introduction à la notion de champ}
% ==================================================

\subsection{Qu'est-ce qu'une interaction ?}

\begin{definition}
On appelle \textbf{interaction} une action réciproque entre deux systèmes.
\end{definition}

Quelques exemples :
\begin{itemize}
    \item interaction gravitationnelle entre la Terre et la Lune ;
    \item interaction électrostatique entre deux corps chargés ;
    \item interaction de contact entre une table et un livre ;
    \item interaction d'un fluide avec une paroi.
\end{itemize}

\begin{remarque}
Dans cette partie du programme, on met surtout l'accent sur :
\begin{itemize}
    \item l'interaction gravitationnelle ;
    \item l'interaction électrostatique.
\end{itemize}
Ces deux interactions peuvent agir à distance.
\end{remarque}

% --------------------------------------------------
\subsection{Charge électrique et interaction électrostatique}
% --------------------------------------------------

\begin{definition}
La \textbf{charge électrique} est une grandeur physique qui caractérise certains objets ou certaines particules.
Elle se note généralement $q$ et s'exprime en coulomb (\si{C}).
\end{definition}

Il existe deux signes de charge :
\begin{itemize}
    \item charge positive ;
    \item charge négative.
\end{itemize}

\begin{propriete}
Deux charges de même signe se repoussent.  
Deux charges de signes opposés s'attirent.
\end{propriete}

\begin{definition}
L'\textbf{influence électrostatique} est le phénomène par lequel un corps chargé modifie la répartition des charges dans un autre corps sans contact.
\end{definition}

\begin{applicationbox}
Quand on approche une règle plastique chargée de petits morceaux de papier, ceux-ci sont d'abord attirés.  
Le papier est globalement neutre, mais la présence de la règle modifie localement la répartition des charges : c'est l'influence électrostatique.
\end{applicationbox}

% --------------------------------------------------
\subsection{Loi de Coulomb}
% --------------------------------------------------

\begin{theoreme}[Loi de Coulomb]
Deux charges ponctuelles $q_A$ et $q_B$, séparées d'une distance $r$, exercent l'une sur l'autre des forces de même direction, de sens opposés, de même valeur :
\[
F = k\,\frac{|q_A q_B|}{r^2}
\]
où $k \approx 9,0\times 10^9\ \si{N.m^2.C^{-2}}$.
\end{theoreme}

La direction de la force est la droite reliant les deux charges.

\begin{itemize}
    \item si $q_A q_B >0$, l'interaction est répulsive ;
    \item si $q_A q_B <0$, l'interaction est attractive.
\end{itemize}

\begin{methodebox}
Pour utiliser la loi de Coulomb :
\begin{enumerate}
    \item repérer les charges et leur signe ;
    \item calculer la distance $r$ ;
    \item utiliser la formule avec les unités SI ;
    \item préciser la direction et le sens de la force ;
    \item ne pas oublier que les deux objets exercent des forces opposées l'un sur l'autre.
\end{enumerate}
\end{methodebox}

\begin{exemple}
Deux charges $q_A=+2,0\times 10^{-6}\ \si{C}$ et $q_B=-3,0\times 10^{-6}\ \si{C}$ sont séparées de \SI{0,20}{m}.  
La valeur de la force vaut :
\[
F = 9,0\times 10^9 \times \frac{2,0\times 10^{-6}\times 3,0\times 10^{-6}}{(0,20)^2}
\]
\[
F = 9,0\times 10^9 \times \frac{6,0\times 10^{-12}}{0,040}
\]
\[
F = 1,35\ \si{N}
\]
Comme les charges sont de signes opposés, la force est attractive.
\end{exemple}

% --------------------------------------------------
\subsection{Gravitation universelle}
% --------------------------------------------------

\begin{theoreme}[Loi de gravitation universelle]
Deux masses $m_A$ et $m_B$, séparées par une distance $r$, exercent l'une sur l'autre des forces attractives de valeur :
\[
F = G\frac{m_A m_B}{r^2}
\]
avec
\[
G=6,67\times 10^{-11}\ \si{N.m^2.kg^{-2}}.
\]
\end{theoreme}

\begin{remarque}
La gravitation agit toujours de manière attractive, contrairement à l'interaction électrostatique qui peut être attractive ou répulsive.
\end{remarque}

\subsection{Analogies entre la loi de Coulomb et la loi de gravitation}

\begin{propriete}
Les deux lois présentent de fortes analogies :
\begin{itemize}
    \item elles décrivent des interactions à distance ;
    \item la valeur de la force est proportionnelle au produit de deux grandeurs caractéristiques :
    \begin{itemize}
        \item les masses pour la gravitation ;
        \item les charges pour l'électrostatique ;
    \end{itemize}
    \item la valeur de la force décroît comme $1/r^2$ ;
    \item la direction est la droite reliant les deux objets.
\end{itemize}
\end{propriete}

\begin{attentionbox}
Différence essentielle :
\begin{itemize}
    \item la gravitation est toujours attractive ;
    \item l'interaction électrostatique peut être attractive ou répulsive.
\end{itemize}
\end{attentionbox}

% --------------------------------------------------
\subsection{De la force au champ}
% --------------------------------------------------

L'idée de champ est extrêmement importante en physique.  
Elle permet de penser qu'un objet modifie l'espace autour de lui, et qu'un autre objet placé dans cette région subit alors une action.

\begin{definition}
Un \textbf{champ} est une grandeur définie en chaque point de l'espace.
\end{definition}

Ici, on introduit :
\begin{itemize}
    \item le champ de gravitation ;
    \item le champ électrostatique.
\end{itemize}

\subsection{Champ de gravitation}

\begin{definition}
Le \textbf{champ de gravitation} créé par une masse $M$ en un point situé à une distance $r$ est défini par :
\[
\vec{g} = \frac{\vec{F}}{m}
\]
où $\vec{F}$ est la force gravitationnelle exercée sur une masse test $m$ placée en ce point.
\end{definition}

Sa valeur est :
\[
g = G\frac{M}{r^2}
\]

Le vecteur $\vec{g}$ est dirigé vers la masse attractive.

\subsection{Champ électrostatique}

\begin{definition}
Le \textbf{champ électrostatique} créé par une charge $Q$ en un point de l'espace est défini par :
\[
\vec{E}=\frac{\vec{F}}{q}
\]
où $\vec{F}$ est la force exercée sur une charge test $q$ placée en ce point.
\end{definition}

Pour une charge ponctuelle $Q$ :
\[
E = k\frac{|Q|}{r^2}
\]

\begin{itemize}
    \item si $Q>0$, le champ est dirigé vers l'extérieur ;
    \item si $Q<0$, le champ est dirigé vers la charge.
\end{itemize}

\begin{methodebox}
Pour passer du champ à la force :
\[
\vec{F}=m\vec{g}
\qquad \text{ou} \qquad
\vec{F}=q\vec{E}
\]
Selon que l'on étudie la gravitation ou l'électrostatique.
\end{methodebox}

% --------------------------------------------------
\subsection{Lignes de champ}
% --------------------------------------------------

\begin{definition}
Une \textbf{ligne de champ} est une courbe telle que le vecteur champ est tangent à cette courbe en chacun de ses points.
\end{definition}

\begin{center}
\begin{tikzpicture}[scale=0.95]
\filldraw[black] (0,0) circle (2pt) node[below] {$+$};
\foreach \a in {0,30,...,330}{
    \draw[->,blue!70!black,thick] (0,0) -- ({2.0*cos(\a)},{2.0*sin(\a)});
}
\node at (0,-2.7) {Lignes de champ d'une charge positive};
\end{tikzpicture}
\end{center}

\begin{remarque}
Plus les lignes de champ sont serrées, plus le champ est intense.
\end{remarque}

\begin{terminalebox}
En Terminale, la notion de champ sera reprise plus profondément, avec davantage de rigueur sur :
\begin{itemize}
    \item la représentation vectorielle ;
    \item les champs uniformes ;
    \item les potentiels et l'énergie potentielle dans certains contextes.
\end{itemize}
\end{terminalebox}

% ==================================================
\section{Description d'un fluide au repos}
% ==================================================

\subsection{Échelle microscopique et échelle macroscopique}

\begin{definition}
Un \textbf{fluide} est un milieu matériel qui peut s'écouler. Les liquides et les gaz sont des fluides.
\end{definition}

On peut décrire un fluide à deux niveaux :

\begin{itemize}
    \item \textbf{microscopique} : on pense aux molécules ou aux atomes qui l'agitent ;
    \item \textbf{macroscopique} : on utilise des grandeurs globales mesurables, comme la pression ou la température.
\end{itemize}

\begin{definition}
Les principales grandeurs macroscopiques décrivant un fluide au repos sont :
\begin{itemize}
    \item la masse volumique $\rho$ ;
    \item la pression $P$ ;
    \item la température $T$.
\end{itemize}
\end{definition}

\subsection{Masse volumique}

\begin{definition}
La \textbf{masse volumique} d'un corps est donnée par :
\[
\rho = \frac{m}{V}
\]
où $m$ est la masse et $V$ le volume.
\end{definition}

Unité :
\[
[\rho]=\si{kg.m^{-3}}
\]

\begin{exemple}
L'eau liquide a une masse volumique voisine de :
\[
\rho_{\text{eau}} \approx \SI{1000}{kg.m^{-3}}
\]
\end{exemple}

\subsection{Interprétation microscopique qualitative}

\begin{remarque}
La pression et la température sont des grandeurs macroscopiques, mais elles sont liées au comportement des particules :

\begin{itemize}
    \item plus les particules frappent fréquemment les parois, plus la pression est grande ;
    \item plus l'agitation microscopique est forte, plus la température est élevée.
\end{itemize}
\end{remarque}

\subsection{Loi de Mariotte}

\begin{theoreme}[Loi de Mariotte]
À température constante, pour une quantité fixe de gaz :
\[
P\times V = \text{constante}
\]
\end{theoreme}

Donc, entre deux états :
\[
P_1V_1=P_2V_2
\]

\begin{attentionbox}
Cette loi concerne un gaz à température constante.  
Elle ne s'applique pas n'importe comment à un liquide ou à un gaz dont la température change.
\end{attentionbox}

\begin{exemple}
Un gaz occupe initialement un volume $V_1=\SI{100}{mL}$ sous une pression $P_1=\SI{1,0e5}{Pa}$.  
On le comprime jusqu'à $V_2=\SI{40}{mL}$. À température constante :
\[
P_2=P_1\frac{V_1}{V_2}
=1,0\times 10^5\times \frac{100}{40}
=2,5\times 10^5\ \si{Pa}
\]
\end{exemple}

% --------------------------------------------------
\subsection{Pression et force pressante}
% --------------------------------------------------

\begin{definition}
La \textbf{pression} est la force pressante exercée par unité de surface :
\[
P=\frac{F}{S}
\]
\end{definition}

Donc :
\[
F=P\times S
\]

Unité :
\[
[P]=\si{Pa}=\si{N.m^{-2}}
\]

\begin{applicationbox}
Si une surface plane de \SI{2,0e-2}{m^2} est soumise à une pression de \SI{1,0e5}{Pa}, alors la force pressante vaut :
\[
F=P\times S = 1,0\times 10^5 \times 2,0\times 10^{-2}=2,0\times 10^3\ \si{N}
\]
\end{applicationbox}

\subsection{Loi fondamentale de la statique des fluides}

\begin{theoreme}
Dans un fluide incompressible au repos :
\[
P_2-P_1=\rho g(z_1-z_2)
\]
où $z$ représente l'altitude.
\end{theoreme}

Si le point 2 est plus bas que le point 1, alors la pression est plus grande en bas.

\begin{propriete}
Dans un liquide au repos, la pression augmente avec la profondeur.
\end{propriete}

\begin{center}
\begin{tikzpicture}[scale=0.95]
\fill[blue!10] (-2,-2) rectangle (2,1.4);
\draw[thick] (-2,1.4)--(2,1.4);
\draw[thick] (-2,-2)--(-2,1.4);
\draw[thick] (2,-2)--(2,1.4);
\filldraw[black] (0,0.8) circle (1.8pt) node[right] {$1$};
\filldraw[black] (0,-1.1) circle (1.8pt) node[right] {$2$};
\draw[dashed] (0,0.8)--(-2.5,0.8);
\draw[dashed] (0,-1.1)--(-2.5,-1.1);
\node[left] at (-2.5,0.8) {$z_1$};
\node[left] at (-2.5,-1.1) {$z_2$};
\node at (0,-2.5) {Dans un liquide au repos, la pression augmente quand l'altitude diminue.};
\end{tikzpicture}
\end{center}

\begin{methodebox}
Pour utiliser la loi de la statique des fluides :
\begin{enumerate}
    \item bien repérer les deux points étudiés ;
    \item choisir le bon axe vertical et les altitudes ;
    \item appliquer la formule en gardant les signes ;
    \item vérifier le sens physique du résultat : plus profond $\Rightarrow$ pression plus grande.
\end{enumerate}
\end{methodebox}

\begin{terminalebox}
En Terminale, tu retrouveras l'idée de variation d'une grandeur dans l'espace, par exemple avec :
\begin{itemize}
    \item les champs ;
    \item les potentiels ;
    \item les bilans d'énergie plus structurés.
\end{itemize}
La statique des fluides prépare déjà à raisonner « localement ».
\end{terminalebox}

% ==================================================
\section{Mouvement d'un système et variation du vecteur vitesse}
% ==================================================

\subsection{Le vecteur variation de vitesse}

\begin{definition}
Entre deux instants voisins $t_1$ et $t_2$, la variation du vecteur vitesse est :
\[
\Delta \vec{v} = \vec{v}_2-\vec{v}_1
\]
\end{definition}

\begin{remarque}
Le vecteur $\Delta \vec{v}$ indique comment la vitesse change :
\begin{itemize}
    \item en valeur ;
    \item en direction ;
    \item en sens.
\end{itemize}
\end{remarque}

\begin{attentionbox}
Même si la valeur de la vitesse reste constante, la variation du vecteur vitesse peut être non nulle si la direction change.  
Exemple : mouvement circulaire.
\end{attentionbox}

% --------------------------------------------------
\subsection{Lien entre variation de vitesse et forces}
% --------------------------------------------------

Le programme de Première introduit une version approchée de la deuxième loi de Newton.

\begin{theoreme}[Relation approchée]
Pour un système modélisé par un point matériel, entre deux instants voisins :
\[
\sum \vec{F} \approx m\frac{\Delta \vec{v}}{\Delta t}
\]
\end{theoreme}

On peut la réécrire :
\[
\Delta \vec{v}\approx \frac{\sum \vec{F}}{m}\Delta t
\]

\begin{propriete}
Cette relation montre que :
\begin{itemize}
    \item plus la somme des forces est grande, plus la vitesse varie rapidement ;
    \item plus la masse est grande, plus il est difficile de modifier le mouvement ;
    \item la variation du vecteur vitesse a la même direction que la somme des forces.
\end{itemize}
\end{propriete}

\begin{remarque}
C'est une manière intuitive d'approcher l'accélération :
\[
\vec{a}\approx \frac{\Delta \vec{v}}{\Delta t}
\]
Ainsi :
\[
\sum \vec{F}\approx m\vec{a}
\]
\end{remarque}

\subsection{Sens physique du rôle de la masse}

La masse mesure en quelque sorte l'inertie du système, c'est-à-dire sa résistance à la modification de son mouvement.

\begin{applicationbox}
À force égale, une balle légère voit sa vitesse varier davantage qu'une boule très massive.
\end{applicationbox}

% --------------------------------------------------
\subsection{Exemples fondamentaux}
% --------------------------------------------------

\subsubsection*{1. Objet sans force résultante}

Si :
\[
\sum \vec{F}=\vec{0}
\]
alors :
\[
\Delta \vec{v}\approx \vec{0}
\]
La vitesse reste approximativement constante.

\subsubsection*{2. Chute verticale sans frottements}

La seule force significative est le poids :
\[
\vec{P}=m\vec{g}
\]
Donc :
\[
\sum \vec{F}=m\vec{g}
\]
La variation de vitesse est dirigée vers le bas.

\subsubsection*{3. Mouvement horizontal avec frottement}

Si un objet glisse horizontalement, la force de frottement est opposée au mouvement.  
La variation du vecteur vitesse est donc opposée à la vitesse : l'objet ralentit.

\subsection{Représentation vectorielle}

\begin{center}
\begin{tikzpicture}[scale=1]
\draw[->,thick] (0,0)--(2.4,0) node[above] {$\vec{v}_1$};
\draw[->,thick] (0,0)--(1.7,1.1) node[above] {$\vec{v}_2$};
\draw[->,red,very thick] (2.4,0)--(1.7,1.1) node[midway,right] {$\Delta \vec{v}$};
\node at (1,-1.0) {Construction graphique de $\Delta\vec{v}=\vec{v}_2-\vec{v}_1$};
\end{tikzpicture}
\end{center}

\begin{methodebox}
Pour construire $\Delta \vec{v}$ :
\begin{enumerate}
    \item placer les deux vecteurs vitesse avec la même origine ;
    \item utiliser la relation $\Delta\vec{v}=\vec{v}_2-\vec{v}_1$ ;
    \item ou bien considérer le vecteur reliant l'extrémité de $\vec{v}_1$ à l'extrémité de $\vec{v}_2$.
\end{enumerate}
\end{methodebox}

% ==================================================
\section{Aspects énergétiques des phénomènes mécaniques}
% ==================================================

\subsection{Pourquoi introduire l'énergie ?}

Décrire le mouvement à partir des forces est très puissant, mais parfois peu pratique.  
Le point de vue énergétique permet d'avoir une vision plus globale :

\begin{itemize}
    \item quelle énergie possède le système ?
    \item comment cette énergie change-t-elle ?
    \item y a-t-il conservation ou dissipation ?
\end{itemize}

\begin{remarque}
Les deux points de vue, dynamique et énergétique, sont complémentaires :
\begin{itemize}
    \item les forces expliquent comment le mouvement change ;
    \item l'énergie permet souvent de calculer plus simplement certains effets globaux.
\end{itemize}
\end{remarque}

% --------------------------------------------------
\subsection{Énergie cinétique}
% --------------------------------------------------

\begin{definition}
L'\textbf{énergie cinétique} d'un système modélisé par un point matériel de masse $m$ et de vitesse $v$ est :
\[
\Ec = \frac{1}{2}mv^2
\]
\end{definition}

Unité :
\[
[\Ec]=\si{J}
\]

\begin{propriete}
L'énergie cinétique :
\begin{itemize}
    \item est toujours positive ou nulle ;
    \item augmente avec la masse ;
    \item augmente très vite avec la vitesse car elle dépend de $v^2$.
\end{itemize}
\end{propriete}

\begin{exemple}
Une balle de masse \SI{0,20}{kg} se déplace à \SI{15}{m.s^{-1}} :
\[
\Ec=\frac12\times 0,20\times 15^2
=22,5\ \si{J}
\]
\end{exemple}

\begin{attentionbox}
Ne pas confondre vitesse et énergie cinétique :
\begin{itemize}
    \item doubler la vitesse ne double pas l'énergie cinétique ;
    \item doubler la vitesse multiplie l'énergie cinétique par 4.
\end{itemize}
\end{attentionbox}

% --------------------------------------------------
\subsection{Travail d'une force}
% --------------------------------------------------

\begin{definition}
Le \textbf{travail} d'une force mesure le transfert d'énergie associé à l'action de cette force lorsque son point d'application se déplace.
\end{definition}

Dans le cas d'une force constante :
\[
W_{A\to B}(\vec{F}) = \vec{F}\cdot \vec{AB}
\]

On peut aussi écrire :
\[
W=F\,AB\,\cos\theta
\]
où $\theta$ est l'angle entre la force et le déplacement.

\begin{propriete}
\begin{itemize}
    \item si le travail est positif, la force apporte de l'énergie au système ;
    \item si le travail est négatif, la force retire de l'énergie au système ;
    \item si le travail est nul, la force ne change pas l'énergie cinétique par ce biais.
\end{itemize}
\end{propriete}

\subsection{Interprétation du signe du travail}

\begin{itemize}
    \item force dans le sens du déplacement : travail positif ;
    \item force opposée au déplacement : travail négatif ;
    \item force perpendiculaire au déplacement : travail nul.
\end{itemize}

\begin{center}
\begin{tikzpicture}[scale=1]
\draw[->,thick] (0,0)--(4,0) node[below] {déplacement};
\draw[->,blue,very thick] (1,0)--(3,0) node[above] {$\vec{F}$};
\node at (2,-1) {Travail positif};

\begin{scope}[xshift=6cm]
\draw[->,thick] (0,0)--(4,0) node[below] {déplacement};
\draw[->,blue,very thick] (3,0)--(1,0) node[above] {$\vec{F}$};
\node at (2,-1) {Travail négatif};
\end{scope}
\end{tikzpicture}
\end{center}

\begin{applicationbox}
Un traîneau est tiré horizontalement sur \SI{8}{m} par une force constante de \SI{50}{N} dans la direction du mouvement.  
Le travail vaut :
\[
W = F\times AB = 50\times 8 = 400\ \si{J}
\]
\end{applicationbox}

% --------------------------------------------------
\subsection{Théorème de l'énergie cinétique}
% --------------------------------------------------

\begin{theoreme}[Théorème de l'énergie cinétique]
La variation d'énergie cinétique d'un système entre deux positions $A$ et $B$ est égale à la somme des travaux des forces qui s'exercent sur lui :
\[
\Ec(B)-\Ec(A)=\sum W_{A\to B}(\vec{F})
\]
\end{theoreme}

\begin{remarque}
C'est un résultat fondamental. Il relie les forces au mouvement, mais via l'énergie.
\end{remarque}

\begin{methodebox}
Pour utiliser le théorème de l'énergie cinétique :
\begin{enumerate}
    \item identifier les forces appliquées ;
    \item calculer le travail de chaque force ;
    \item sommer les travaux ;
    \item écrire la variation d'énergie cinétique ;
    \item résoudre pour la grandeur demandée.
\end{enumerate}
\end{methodebox}

\subsection{Exemple complet}

Un objet de masse \SI{2,0}{kg} se déplace à \SI{3,0}{m.s^{-1}} puis atteint \SI{7,0}{m.s^{-1}}.  
La variation d'énergie cinétique vaut :
\[
\Delta \Ec=\frac12 m(v_B^2-v_A^2)
\]
\[
\Delta \Ec=\frac12 \times 2,0 \times (7,0^2-3,0^2)
\]
\[
\Delta \Ec=1\times (49-9)=40\ \si{J}
\]
Donc la somme des travaux des forces vaut \SI{40}{J}.

% --------------------------------------------------
\subsection{Forces conservatives et énergie potentielle}
% --------------------------------------------------

\begin{definition}
Une \textbf{force conservative} est une force pour laquelle le travail entre deux points ne dépend pas du chemin suivi, mais seulement de la position initiale et de la position finale.
\end{definition}

Exemple majeur du programme : le poids dans le champ de pesanteur terrestre, près de la surface de la Terre.

\begin{definition}
À une force conservative, on peut associer une \textbf{énergie potentielle}.
\end{definition}

\subsection{Énergie potentielle de pesanteur}

\begin{theoreme}
Au voisinage de la surface de la Terre, l'énergie potentielle de pesanteur d'un système de masse $m$ à l'altitude $z$ s'écrit :
\[
\Ep = mgz + \text{constante}
\]
\end{theoreme}

En choisissant le niveau de référence où $\Ep=0$, on peut écrire simplement :
\[
\Ep = mgz
\]

\begin{remarque}
L'énergie potentielle dépend du choix de l'origine des altitudes.  
Ce qui compte physiquement, c'est sa variation.
\end{remarque}

\subsection{Démonstration de la variation d'énergie potentielle de pesanteur}

Considérons un déplacement vertical d'un point $A$ vers un point $B$.

Le poids vaut :
\[
\vec{P}=m\vec{g}
\]

Le travail du poids entre $A$ et $B$ s'écrit :
\[
W_{A\to B}(\vec{P}) = mg(z_A-z_B)
\]

On définit l'énergie potentielle de sorte que :
\[
W_{A\to B}(\vec{P}) = \Ep(A)-\Ep(B)
\]

On obtient alors :
\[
\Ep(B)-\Ep(A)=mg(z_B-z_A)
\]

Une expression compatible est :
\[
\Ep = mgz
\]

\begin{attentionbox}
Quand l'altitude augmente, l'énergie potentielle augmente.  
Quand un objet descend, son énergie potentielle diminue.
\end{attentionbox}

% --------------------------------------------------
\subsection{Forces non conservatives : frottements}
% --------------------------------------------------

\begin{definition}
Une \textbf{force non conservative} dissipe une partie de l'énergie mécanique, souvent sous forme thermique.
\end{definition}

Exemple : les frottements.

Si une force de frottement d'intensité constante $f$ s'oppose au mouvement rectiligne sur une distance $d$, alors :
\[
W(\vec{f})=-fd
\]

\begin{remarque}
Le travail des frottements est négatif : ils retirent de l'énergie mécanique au système.
\end{remarque}

% --------------------------------------------------
\subsection{Énergie mécanique}
% --------------------------------------------------

\begin{definition}
L'\textbf{énergie mécanique} d'un système est la somme de son énergie cinétique et de son énergie potentielle :
\[
\Em = \Ec + \Ep
\]
\end{definition}

\subsection{Conservation de l'énergie mécanique}

\begin{theoreme}
Si seules des forces conservatives travaillent, l'énergie mécanique se conserve :
\[
\Em = \text{constante}
\]
\end{theoreme}

\begin{propriete}
En l'absence de frottements :
\begin{itemize}
    \item une diminution de l'énergie potentielle s'accompagne d'une augmentation de l'énergie cinétique ;
    \item et réciproquement.
\end{itemize}
\end{propriete}

\begin{exemple}
En chute libre sans frottements :
\[
\Ec + \Ep = \text{constante}
\]
Quand l'objet descend, $\Ep$ diminue et $\Ec$ augmente.
\end{exemple}

\subsection{Non-conservation de l'énergie mécanique}

Si des forces non conservatives travaillent :
\[
\Delta \Em = W(\text{forces non conservatives})
\]

En particulier, avec des frottements :
\[
\Delta \Em <0
\]

\begin{remarque}
L'énergie mécanique n'est pas « détruite » au sens absolu : elle est transformée vers d'autres formes, souvent de l'énergie interne ou thermique.
\end{remarque}

\begin{terminalebox}
En Terminale, les bilans énergétiques deviendront plus fins.  
Tu verras que le raisonnement par conservation de l'énergie est extrêmement puissant, notamment pour :
\begin{itemize}
    \item les oscillateurs ;
    \item les circuits électriques ;
    \item certaines transformations thermodynamiques ;
    \item les mouvements dans des champs.
\end{itemize}
\end{terminalebox}

% ==================================================
\section{Aspects énergétiques des phénomènes électriques}
% ==================================================

\subsection{Porteurs de charge et intensité du courant}

\begin{definition}
Le courant électrique correspond à un déplacement ordonné de porteurs de charge.
\end{definition}

\begin{definition}
L'intensité du courant est le débit de charges :
\[
I=\frac{\Delta Q}{\Delta t}
\]
\end{definition}

Unité :
\[
[I]=\si{A}
\]

\begin{remarque}
Cette formule signifie qu'une intensité grande correspond à une grande quantité de charge traversant une section en peu de temps.
\end{remarque}

\begin{exemple}
Si \SI{12}{C} traversent une section en \SI{3}{s}, alors :
\[
I=\frac{12}{3}=4,0\ \si{A}
\]
\end{exemple}

% --------------------------------------------------
\subsection{Source réelle de tension}
% --------------------------------------------------

\begin{definition}
Une source réelle de tension continue peut être modélisée par l'association en série :
\begin{itemize}
    \item d'une source idéale de tension ;
    \item d'une résistance interne.
\end{itemize}
\end{definition}

\begin{remarque}
Cela permet d'expliquer pourquoi la tension délivrée par une pile réelle dépend du courant qu'elle fournit.
\end{remarque}

\begin{attentionbox}
Une pile réelle n'est pas un générateur parfait : une partie de l'énergie est dissipée à l'intérieur même de la source.
\end{attentionbox}

% --------------------------------------------------
\subsection{Puissance, énergie et effet Joule}
% --------------------------------------------------

\begin{definition}
La \textbf{puissance} est le débit d'énergie :
\[
P=\frac{E}{\Delta t}
\]
\end{definition}

Donc :
\[
E=P\Delta t
\]

Dans un dipôle électrique :
\[
P=UI
\]

Pour un dipôle ohmique :
\[
P=RI^2=\frac{U^2}{R}
\]

\begin{definition}
L'\textbf{effet Joule} correspond à la dissipation d'énergie électrique sous forme thermique dans un conducteur traversé par un courant.
\end{definition}

\begin{exemple}
Une résistance de \SI{10}{\ohm} est traversée par un courant de \SI{2,0}{A}.  
La puissance dissipée est :
\[
P=RI^2=10\times 2,0^2=40\ \si{W}
\]
\end{exemple}

% --------------------------------------------------
\subsection{Rendement}
% --------------------------------------------------

\begin{definition}
Le \textbf{rendement} d'un convertisseur est :
\[
\eta = \frac{P_{\text{utile}}}{P_{\text{reçue}}}
\]
ou de manière équivalente :
\[
\eta=\frac{E_{\text{utile}}}{E_{\text{reçue}}}
\]
\end{definition}

\begin{propriete}
On a toujours :
\[
0\leq \eta \leq 1
\]
souvent exprimé en pourcentage.
\end{propriete}

\begin{applicationbox}
Un moteur reçoit une puissance électrique de \SI{120}{W} et fournit une puissance mécanique utile de \SI{90}{W}.  
Son rendement vaut :
\[
\eta=\frac{90}{120}=0,75
\]
soit \SI{75}{\percent}.
\end{applicationbox}

% ==================================================
\section{Grandes méthodes de résolution}
% ==================================================

\subsection{Méthode 1 : résoudre un exercice de forces et mouvement}

\begin{methodebox}
\begin{enumerate}
    \item Identifier le système étudié.
    \item Choisir le référentiel.
    \item Faire le bilan des forces.
    \item Représenter les vecteurs.
    \item Utiliser la relation :
    \[
    \sum \vec{F}\approx m\frac{\Delta \vec{v}}{\Delta t}
    \]
    \item Interpréter physiquement le résultat.
\end{enumerate}
\end{methodebox}

\subsection{Méthode 2 : résoudre un exercice avec l'énergie}

\begin{methodebox}
\begin{enumerate}
    \item Définir les états initial et final.
    \item Calculer $\Ec$, puis éventuellement $\Ep$.
    \item Écrire soit :
    \[
    \Delta \Ec = \sum W(\vec{F})
    \]
    soit :
    \[
    \Delta \Em = W(\text{forces non conservatives})
    \]
    \item Exploiter la relation adaptée à la situation.
    \item Vérifier la cohérence physique.
\end{enumerate}
\end{methodebox}

\subsection{Méthode 3 : résoudre un exercice de statique des fluides}

\begin{methodebox}
\begin{enumerate}
    \item Identifier les deux points.
    \item Repérer les altitudes.
    \item Utiliser :
    \[
    P_2-P_1=\rho g(z_1-z_2)
    \]
    \item Vérifier que le point le plus bas a bien la pression la plus grande.
\end{enumerate}
\end{methodebox}

% ==================================================
\section{Tableau de synthèse}
% ==================================================

\begin{center}
\renewcommand{\arraystretch}{1.5}
\begin{tabularx}{\textwidth}{|>{\raggedright\arraybackslash}p{4.2cm}|X|X|}
\hline
\textbf{Notion} & \textbf{Expression} & \textbf{Sens physique} \\
\hline
Loi de Coulomb & $F=k\dfrac{|q_Aq_B|}{r^2}$ & Interaction électrostatique à distance \\
\hline
Gravitation & $F=G\dfrac{m_Am_B}{r^2}$ & Attraction gravitationnelle à distance \\
\hline
Champ de gravitation & $\vec{g}=\dfrac{\vec{F}}{m}$ & Action gravitationnelle par unité de masse \\
\hline
Champ électrostatique & $\vec{E}=\dfrac{\vec{F}}{q}$ & Action électrique par unité de charge \\
\hline
Masse volumique & $\rho=\dfrac{m}{V}$ & Quantité de masse par unité de volume \\
\hline
Loi de Mariotte & $P_1V_1=P_2V_2$ & Compression/détente isotherme d'un gaz \\
\hline
Force pressante & $F=PS$ & Action d'un fluide sur une surface \\
\hline
Statique des fluides & $P_2-P_1=\rho g(z_1-z_2)$ & Variation de pression avec l'altitude \\
\hline
Variation de vitesse & $\sum \vec{F}\approx m\dfrac{\Delta \vec{v}}{\Delta t}$ & Lien entre forces et changement de mouvement \\
\hline
Énergie cinétique & $\Ec=\dfrac12 mv^2$ & Énergie liée au mouvement \\
\hline
Travail d'une force constante & $W=\vec{F}\cdot \vec{AB}$ & Transfert d'énergie par une force \\
\hline
Théorème de l'énergie cinétique & $\Delta \Ec = \sum W$ & Bilan énergétique du mouvement \\
\hline
Énergie potentielle de pesanteur & $\Ep=mgz$ & Énergie liée à la position verticale \\
\hline
Énergie mécanique & $\Em=\Ec+\Ep$ & Énergie globale mécanique \\
\hline
Courant continu & $I=\dfrac{\Delta Q}{\Delta t}$ & Débit de charges \\
\hline
Puissance électrique & $P=UI$ & Débit d'énergie électrique \\
\hline
Rendement & $\eta=\dfrac{P_{\text{utile}}}{P_{\text{reçue}}}$ & Efficacité d'un convertisseur \\
\hline
\end{tabularx}
\end{center}

% ==================================================
\section{Vrai/Faux commentés}
% ==================================================

\begin{enumerate}[label=\textbf{\arabic*.}]
    \item \textbf{« Si un objet se déplace, alors une force agit forcément dans le sens du mouvement. »}  
    Faux. Un objet peut se déplacer à vitesse constante sans force résultante non nulle.

    \item \textbf{« Deux charges de même signe s'attirent. »}  
    Faux. Elles se repoussent.

    \item \textbf{« La gravitation et l'interaction électrostatique décroissent toutes deux en $1/r^2$. »}  
    Vrai.

    \item \textbf{« Dans un liquide au repos, la pression diminue quand on descend. »}  
    Faux. Elle augmente avec la profondeur.

    \item \textbf{« Si la somme des forces est nulle, alors la variation du vecteur vitesse est nulle. »}  
    Vrai, dans le cadre du modèle étudié.

    \item \textbf{« Le travail d'une force perpendiculaire au déplacement est nul. »}  
    Vrai.

    \item \textbf{« Une force de frottement peut augmenter l'énergie mécanique. »}  
    Faux dans les situations usuelles du programme : elle la dissipe.

    \item \textbf{« L'énergie cinétique dépend de la vitesse au carré. »}  
    Vrai.

    \item \textbf{« Une pile réelle peut être modélisée par une source idéale seule. »}  
    Faux. Il faut tenir compte de sa résistance interne.

    \item \textbf{« Le rendement peut dépasser 100\,\%. »}  
    Faux.
\end{enumerate}

% ==================================================
\section{Exercices d'application corrigés}
% ==================================================

\subsection*{Exercice 1 -- Interaction électrostatique}

Deux charges $q_A=+1,0\times10^{-6}\ \si{C}$ et $q_B=+2,0\times10^{-6}\ \si{C}$ sont séparées de \SI{0,30}{m}.  
Calculer la valeur de la force électrostatique.

\paragraph{Corrigé.}
\[
F=k\frac{|q_Aq_B|}{r^2}
\]
\[
F=9,0\times10^9\times \frac{(1,0\times10^{-6})(2,0\times10^{-6})}{0,30^2}
\]
\[
F=9,0\times10^9\times \frac{2,0\times10^{-12}}{0,09}
\]
\[
F=0,20\ \si{N}
\]
Les charges ont le même signe : la force est répulsive.

\subsection*{Exercice 2 -- Pression dans un liquide}

On considère de l'eau de masse volumique $\rho=\SI{1000}{kg.m^{-3}}$.  
Deux points sont séparés verticalement de \SI{2,0}{m}, le second étant plus bas.  
Calculer la différence de pression.

\paragraph{Corrigé.}
\[
P_2-P_1=\rho g(z_1-z_2)
\]
Ici, $z_1-z_2=2,0$ m, donc :
\[
P_2-P_1=1000\times 9,81\times 2,0
\]
\[
P_2-P_1=1,96\times10^4\ \si{Pa}
\]

\subsection*{Exercice 3 -- Théorème de l'énergie cinétique}

Un objet de masse \SI{1,5}{kg} passe de \SI{2,0}{m.s^{-1}} à \SI{6,0}{m.s^{-1}}.  
Calculer la somme des travaux des forces.

\paragraph{Corrigé.}
\[
\Delta \Ec = \frac12 m(v_B^2-v_A^2)
\]
\[
\Delta \Ec=\frac12\times1,5\times(36-4)
\]
\[
\Delta \Ec=0,75\times 32 = 24\ \si{J}
\]
La somme des travaux des forces vaut donc \SI{24}{J}.

\subsection*{Exercice 4 -- Énergie mécanique}

Une bille de masse \SI{0,20}{kg} est lâchée sans vitesse initiale d'une hauteur \SI{1,5}{m}.  
On néglige les frottements.  
Calculer sa vitesse juste avant l'impact avec le sol.

\paragraph{Corrigé.}
Conservation de l'énergie mécanique :
\[
\Ec_i+\Ep_i=\Ec_f+\Ep_f
\]
Initialement :
\[
\Ec_i=0,\qquad \Ep_i=mgh
\]
Finalement au sol :
\[
\Ep_f=0,\qquad \Ec_f=\frac12 mv^2
\]
Donc :
\[
mgh=\frac12 mv^2
\]
\[
gh=\frac12 v^2
\]
\[
v^2=2gh
\]
\[
v=\sqrt{2\times 9,81\times 1,5}
\approx \SI{5,4}{m.s^{-1}}
\]

% ==================================================
\section{QCM de compréhension}
% ==================================================

\subsection*{QCM}

Pour chaque question, une ou plusieurs réponses peuvent être justes.

\begin{enumerate}[label=\textbf{Q\arabic*.}]
    \item La force gravitationnelle :
    \begin{enumerate}[label=\alph*)]
        \item peut être attractive ou répulsive ;
        \item est toujours attractive ;
        \item décroît comme $1/r$ ;
        \item décroît comme $1/r^2$.
    \end{enumerate}

    \item Dans un liquide au repos :
    \begin{enumerate}[label=\alph*)]
        \item la pression est la même partout ;
        \item la pression augmente avec la profondeur ;
        \item la pression dépend de la masse volumique ;
        \item la pression dépend uniquement du volume du récipient.
    \end{enumerate}

    \item Si le travail d'une force est négatif :
    \begin{enumerate}[label=\alph*)]
        \item la force favorise le mouvement ;
        \item la force retire de l'énergie au système ;
        \item la force est nécessairement perpendiculaire au déplacement ;
        \item la force s'oppose globalement au déplacement.
    \end{enumerate}

    \item L'énergie mécanique :
    \begin{enumerate}[label=\alph*)]
        \item est toujours conservée ;
        \item est la somme de l'énergie cinétique et de l'énergie potentielle ;
        \item peut diminuer à cause des frottements ;
        \item est une force.
    \end{enumerate}

    \item Le rendement :
    \begin{enumerate}[label=\alph*)]
        \item peut être supérieur à 1 ;
        \item mesure l'efficacité d'un convertisseur ;
        \item vaut le rapport de l'énergie utile à l'énergie reçue ;
        \item s'exprime souvent en pourcentage.
    \end{enumerate}
\end{enumerate}

\subsection*{Correction du QCM}

\begin{itemize}
    \item Q1 : b, d
    \item Q2 : b, c
    \item Q3 : b, d
    \item Q4 : b, c
    \item Q5 : b, c, d
\end{itemize}

% ==================================================
\section{Ce qu'il faut absolument savoir faire}
% ==================================================

\begin{objectifbox}
À maîtriser parfaitement :
\begin{itemize}
    \item calculer une force gravitationnelle ou électrostatique ;
    \item comparer gravitation et électrostatique ;
    \item utiliser la notion de champ et interpréter des lignes de champ ;
    \item utiliser $P_1V_1=P_2V_2$ ;
    \item utiliser $F=PS$ ;
    \item utiliser la loi de statique des fluides ;
    \item relier somme des forces et variation du vecteur vitesse ;
    \item calculer une énergie cinétique ;
    \item calculer le travail d'une force constante ;
    \item utiliser le théorème de l'énergie cinétique ;
    \item utiliser l'énergie potentielle de pesanteur ;
    \item raisonner sur la conservation ou non de l'énergie mécanique ;
    \item relier intensité et débit de charges ;
    \item calculer une puissance, une énergie et un rendement.
\end{itemize}
\end{objectifbox}

% ==================================================
\section*{Conclusion}
% ==================================================

Le thème \textbf{Mouvement et interactions} est l'un des plus importants de la physique de Première.  
Il fournit les bases d'une pensée physique rigoureuse :

\begin{itemize}
    \item décrire un système ;
    \item identifier les actions qu'il subit ;
    \item prévoir l'évolution de son mouvement ;
    \item interpréter cette évolution en termes de forces ou d'énergie.
\end{itemize}

C'est aussi un thème profondément unificateur :
\begin{itemize}
    \item la gravitation relie la chute d'un objet et le mouvement des planètes ;
    \item les champs relient action à distance et représentation spatiale ;
    \item l'énergie relie mouvement, transformation et efficacité.
\end{itemize}

\begin{terminalebox}
Si ce chapitre est bien compris, la Terminale sera beaucoup plus fluide, car tu auras déjà les bons réflexes :
\begin{itemize}
    \item toujours partir d'un système bien défini ;
    \item faire un bilan des forces ;
    \item savoir choisir entre approche dynamique et approche énergétique ;
    \item vérifier la cohérence physique des résultats.
\end{itemize}
\end{terminalebox}

\end{document}